\section{WBA for oscillatory signals}
\label{sec:WBAforDualFrequencySignals}
%\subsection{Analysis of the problem}
In section \ref{sec:SignalsWithTwoFrequencies} we encounter a problem with the frequency analysis for our artificial oscillatory signals with certain combinations of $\nu_{1,2}$. 
The claim is, that the analysis works if and only if the signal given by $\Delta\phi_n\,\tecxt\text{mod}\,1$ as a function of $\phi_n$ is in one piece, i.e. shows no large gaps in $\Delta\phi_n$ relative to the full extent. 
%The motivation for the following is, that we compute the frequency as the WBA of the $\Delta\phi_n$. 
Here $\Delta\phi_n:=\phi_{n+1}-\phi_n$ are the angle differences, and the $\phi_n$ come from the transformation of $q_n$ and $p_n$ to polar coordinates. 
To understand this, it is helpful to first leave out the modular arithmetic, which leads to the signals shown in fig. \ref{figure:OscPeriodic2fA05PhiDiffNoMod}. 
Here $\Delta\phi_n$ is shown for $\nu_1=\sqrt{2}-1$, and $\nu_2=\frac{1}{8}$ in blue and for swapped frequencies in orange. 
Both signals are created with an amplitude of $A=0.5$ from eqn. (\ref{eqn:TwoFrequencySignal}).
%, that are otherwise identical to the example from \ref{sec:SignalsWithTwoFrequencies}. 
\begin{figure}[h!]
\centering
\includegraphics[width=\linewidth]{pictures/OscPeriodic2fA05PhiDiffNoMod.png}
\caption[Angle differences for two dual frequency signals before modulo operation]{Angle differences $\Delta\phi_n$ as a function of $\phi_n$. 
The angles $\phi_n$ are obtained from eqn. (\ref{eqn:TwoFrequencySignal}) with an amplitude of $A=0.5$.
The frequencies are $\nu_1=\sqrt{2}-1$ and $\nu_2=\frac{1}{8}$ for the blue points, and are swapped for the orange points.}
\label{figure:OscPeriodic2fA05PhiDiffNoMod}
\end{figure}
\\
Notice how the resulting signals seem to be torn into two pieces, which is problematic for the frequency analysis. 
Applying the WBA from eqn. (\ref{eqn:WBAdefinition}) to $\Delta\phi_n$ results in a value close to zero. 
This is because the clockwise and counter clockwise rotation are indistinguishable, so the angle differences average out. 
Therefore, it is necessary to apply some transformation to get a meaningful result. 
As the gaps are a consequence of the discontinuity of the polar angle, the natural solution seems to be to remove those gaps. 
%The second problem is then purely a question of how to accomplish this. 
\\
Let us introduce the parameterized transformation
\begin{align*}
\Delta\phi_n&\longrightarrow\klam\left[ {\kla\left( {\Delta\phi_n+\alpha} \right)\,\tecxt\text{mod}\,1} \right]-\alpha\numberthis\label{eqn:DeltaPhiModAlpha}
\end{align*}
for $\alpha\in [0,1)$ as a generalization of eqn. (\ref{eqn:pnFromPhiDiff}). 
The case $\alpha=0$ for the signals from fig. \ref{figure:OscPeriodic2fA05PhiDiffNoMod} is already shown in fig. \ref{figure:OscPeriodic2fA05PhiDiff}. 
The frequency pair is chosen in such a way, that the transformation only works for one of the two signals. 
In order to show, that eqn. (\ref{eqn:DeltaPhiModAlpha}) is not a general solution for any $\alpha$, we apply it to signals for the complete $(\nu_1,\nu_2)$-frequency space, which is shown in fig. \ref{figure:OscPeriodic2fA05GridAlpha}. 
The evenly spaced grid is created with a step size of $\delta:=\frac{\varphi_\tecxt\text{golden}}{180}$ and an initial value $\nu_{1,2}=\frac{\delta}{2}$ for the frequencies, resulting a size of about $300\times 300$ points.
For each pair, the accuracy of the frequency analysis is represented as a color value, where the signals are again created using eqn. (\ref{eqn:TwoFrequencySignal}) with $A=0.5$ and a length of $N=1024$.
The angle differences are transformed according to eqn. (\ref{eqn:DeltaPhiModAlpha}) for four different $\alpha$.
This gives rise to a critical region, where the calculated frequency is incorrect.
\begin{figure}[h!]
\centering
\includegraphics[width=\linewidth]{pictures/OscPeriodic2fA05GridAlphaWBA.png}
%\begin{tabular}{ll}
%\includegraphics[width=.5\linewidth]{pictures/OscPeriodic2fA05GridN10Alpha000WBA.png}&
%\includegraphics[width=.5\linewidth]{pictures/OscPeriodic2fA05GridN10Alpha015WBA.png}\\
%\includegraphics[width=.5\linewidth]{pictures/OscPeriodic2fA05GridN10Alpha035WBA.png}&
%\includegraphics[width=.5\linewidth]{pictures/OscPeriodic2fA05GridN10Alpha050WBA.png}\\
%\end{tabular}
\caption[Critical region in the colormap of the convergence of WBA for dual frequency signals]{The color value corresponds to $|\nu_1-\nu_N|$ for signals with two frequencies $\nu_{1,2}$ and a relative amplitude of $A=0.5$ from eqn. (\ref{eqn:TwoFrequencySignal}). 
The signal length for all plots is $N=1024$, and the $\nu_N$ are calculated as the WBA of the angle differences from eqn. (\ref{eqn:DeltaPhiModAlpha}). 
Notice how the critical region moves along the line given by $\nu_2=\nu_1$, as $\alpha$ is varied.}
\label{figure:OscPeriodic2fA05GridAlpha}
\end{figure}
\\
This critical region is a consequence of the modulo operation in eqn. (\ref{eqn:DeltaPhiModAlpha}), that is necessary to get a value different from zero for the frequency, but splits some signals into two parts in the way described above.
The problem persists for various values of $\alpha$, which demonstrates, that eqn. (\ref{eqn:DeltaPhiModAlpha}) can not be used as a general solution.
\\
%It turns out that this is precisely the criterion, for which the transformation described in (\ref{eqn:pnFromPhiDiff}) fails to retrieve the correct frequency with WBA. 
%\FloatBarrier
%\subsection{Solutions to the problem}
There may be several approaches to solve this problem.
Our algorithm follows the basic principle presented in \cite{DasSaiSanYor2019}, but does not utilize their embedding of the torus.
The underlying assumption is, that the angle differences are only defined up to an integer.
The angles are $\phi_n\in [-0.5,0.5)$, and hence the differences are $\Delta\phi_n\in (-1, 1)$.
Because we compute our frequency $\nu\in[0,1)$ as the WBA of the $\Delta\phi_n$, it should be sufficient for the angle differences to also lie within an interval of length one.
\\
More formally, let these new angle differences be defined as
\begin{align*}
\hat{\Delta}\phi_n:=\Delta\phi_n -k_n\numberthis\label{eqn:HatDeltaPhi}
\end{align*}
for integers $k_n\in\mathbb{Z}$, which we call the compartment of each $\Delta\phi_n$, such that all $\hat{\Delta}\phi_n$ lie within an interval of length one.
Equivalently, our goal is to find values for $k_n$, such that the $\hat{\Delta}\phi_n$ all lie within the same compartment.
As a side note, the visual impression of a signal being split in two parts as in fig. \ref{figure:OscPeriodic2fA05PhiDiffNoMod}, can now be described by $k_n\neq k_m$ for some $n,m$.
\\
Define the angle difference $\Delta\phi_m$ for some index $m$ as our point of reference, i.e. choose $k_m=0$.
Now define the differences
\begin{align*}
\delta_{nm}&:=\Delta\phi_n - \Delta\phi_m&\tecxt\text{and}&&\hat{\delta}_{nm}&:=\hat{\Delta}\phi_n - \hat{\Delta}\phi_m, \numberthis\label{eqn:DeltaNM}
\end{align*}
which are related by
\begin{align*}
\delta_{nm}&=\hat{\delta}_{nm}+k_n.
\end{align*}
Now we only need to find the $k_n$, for which the $\hat{\delta}_{nm}$ are as close to zero as possible.
Our algorithm computes the $k_n$ based on the rules
\begin{align*}
k_n&=\begin{cases}
+1&\tecxt\text{if}\ \delta_{nm}>1/2 \\
-1&\tecxt\text{if}\ \delta_{nm}<-1/2 \\
\ 0 &\tecxt\text{else}
\end{cases} .\numberthis\label{eqn:RulesForKN}
\end{align*}
The simplicity may also coincide with a potential loss of generality, since we encounter a problem for the type of orbit presented in the next section.
However, it is unclear, whether the algorithm in \cite{DasSaiSanYor2019} would be able to resolve this either.
%One simple solution is based on the two signals from fig. \ref{figure:OscPeriodic2fA05PhiDiffNoMod}. 
%As we have seen, taking the modulus of $\Delta\phi_n$ results in a correct and an incorrect signal. 
%We want a systematic way of shifting only the bottom part of $\Delta\phi_n$ by 1. 
%For this it turns out to be sufficient, to check whether any parts of the signal lie in both $I_1=[0,\frac{1}{4})$ and $I_2=(\frac{3}{4},1]$. 
%If this condition is met, we transform $\Delta\phi_n\longrightarrow\Delta\phi_n + 1$, for all $\Delta\phi_n<\frac{1}{2}$. 
%This removes the critical region in the case $\alpha=0$ from fig. \ref{figure:OscPeriodic2fA05GridAlpha}. 
%It is even possible to generalize this approach for all $\alpha$. 
%The condition becomes based on the intervals 
%\begin{align*}
%I_1&:=\left[-\alpha,\frac{1}{4}-\alpha\right)&\tecxt\text{and}&&I_2&:=\left(\frac{3}{4}-\alpha,1-\alpha\right],
%\end{align*}
%while the transformation now reads
%\begin{align*}
%\Delta\phi_n&\longrightarrow\Delta\phi_n + 1&\forall n:\ \Delta\phi_n<\frac{1}{2}-\alpha.
%\end{align*}
%For convenience, a value of $\alpha=\frac{1}{2}$ was chosen, as that slightly simplified the implementation.
%\\
%A different approach could be to apply (\ref{eqn:DeltaPhiModAlpha}) to $\Delta\phi_n$ with a value of $\alpha$, that may depend on the specific signal. 
%A candidate for such a shift is given by $\alpha=-\nu$, which does not seem particularly helpful, as we have to transform the signal to retrieve the frequency in the first place. 
%However, a very rough approximation of said frequency would already suffice, as the value of $\alpha$ can generally be chosen from a certain range.
\section{Failure of WBA for a specific type of orbit}
\label{sec:FailureWBASpecificOrbit}
In this section we want to demonstrate the failure of the WBA for some of the four dimensional orbits from section \ref{sec:FrequencySpace}.
They appear to be particularly common near edges of the frequency space beyond the $3:1:1$ and $-1:2:0$ resonances, as well as on certain resonance lines.
Panels (a) and (b) of fig. \ref{figure:Torus4dTransformDeltaPhi} show the projection of a transformed orbit, that is used to compute the second frequency, for two examples.
Respectively, panels (c) and (d) show a plot of the angle differences $\hat{\Delta}\phi_n$, obtained from the algorithm in eqn. (\ref{eqn:RulesForKN}), as a function of the scaled polar angles $\phi_n$ from eqn. (\ref{eqn:MapArctan2}). 
\begin{figure}[h!]
\centering
\includegraphics[width=\linewidth]{pictures/Torus4dTransform_DeltaPhi.png}
\caption[Comparison of the projections and angle differences for computing $\nu_2$ with WBA in two examples]{Shown in (a) and (b) are two examples of the projection of a transformed orbit, which is used for computing the second frequency in either case.
Shown in (c) and (d), corresponding to (a) and (b) respectively, is a plot of the angle differences $\hat{\Delta}\phi_n$.
These are obtained from the algorithm in eqn. (\ref{eqn:RulesForKN}), as a function of the scaled polar angles $\phi_n$ from eqn. (\ref{eqn:MapArctan2}).
The projection in (b) shows no central hole, and the respective $\hat{\Delta}\phi_n$ in (d) span over an interval of length one.
While the two seem to be connected, the more precise statement is, that the frequency analysis with WBA fails, when the second criterion is met.}%WBA:  0.14756, Naff: 0.14756 and WBA:  0.08614, Naff: 0.07592
\label{figure:Torus4dTransformDeltaPhi}
\end{figure}
\\
For the example in (a), the projection shows a hole, and the $\hat{\Delta}\phi_n$ only span over a small interval.
Here the frequency analysis with WBA works, i.e. gives the same result as Naff.
In (b), the central hole is missing, and the $\hat{\Delta}\phi_n$ fill out the largest possible interval, since the length is limited to 1 by construction.
Our claim is, that the WBA fails whenever this second criterion is met.
In all the considered examples, the first criterion is also met, but a more extensive analysis would be required to back up this hypothetical connection.
We could not reconstruct the ``embedding continuation method'' presented in \cite{DasSaiSanYor2019}, so it may provide a solution to this problem.
However, none of the examples considered in this citation are -- for the lack of a better word -- quite as messy as the one from fig. \ref{figure:Torus4dTransformDeltaPhi} (b).